% =====================================================================
%  Experimental Protocols for the Validation of Circuit-Based
%  Disease Diagnosis via the Cause-for-Concern Calculus
%
%  Designs five experiments against public data, specifies the
%  CFC syntax extensions each requires, states success and failure
%  criteria, and gives the complete CFC programs.
% =====================================================================
\documentclass[11pt,a4paper]{article}

\usepackage[utf8]{inputenc}
\usepackage[T1]{fontenc}
\usepackage{lmodern}
\usepackage[a4paper,margin=2.4cm]{geometry}
\usepackage{amsmath,amssymb,amsthm,mathtools}
\usepackage{bm}
\usepackage{booktabs}
\usepackage{array}
\usepackage{enumitem}
\usepackage{microtype}
\usepackage{xcolor}
\usepackage{listings}
\usepackage{algorithm}
\usepackage{algpseudocode}
\usepackage[round,sort&compress]{natbib}
\usepackage[colorlinks=true,linkcolor=blue!45!black,
            citecolor=blue!45!black,urlcolor=blue!45!black]{hyperref}
\usepackage{cleveref}
\usepackage{longtable}
\usepackage{multirow}

% ---- Theorem environments ----
\newtheoremstyle{expthm}{\topsep}{\topsep}{\itshape}{}{\bfseries}{.}{.5em}{}
\theoremstyle{expthm}
\newtheorem{theorem}{Theorem}[section]
\newtheorem{lemma}[theorem]{Lemma}
\newtheorem{proposition}[theorem]{Proposition}
\newtheorem{corollary}[theorem]{Corollary}
\theoremstyle{definition}
\newtheorem{definition}[theorem]{Definition}
\newtheorem{assumption}{Assumption}
\newtheorem{principle}[theorem]{Principle}
\newtheorem{protocol}[theorem]{Protocol}
\newtheorem{criterion}[theorem]{Criterion}
\newtheorem{example}[theorem]{Example}
\theoremstyle{remark}
\newtheorem{remark}[theorem]{Remark}

% ---- .cfc listing style ----
\definecolor{cfckw}{RGB}{30,64,175}
\definecolor{cfccom}{RGB}{22,101,52}
\definecolor{cfcnum}{RGB}{180,83,9}
\definecolor{cfcstr}{RGB}{159,18,57}
\definecolor{cfcbg}{RGB}{250,250,250}
\lstdefinelanguage{cfc}{
  morekeywords={floor,patient,probe,panel,measure,localize,cell,gap,
    close,toward,coherent,using,basis,until,resolved,contested,decline,
    concern,when,emit,assert,yield,let,in,as,with,import,module,export,
    monitor,every,observe,against,reference,circuit,from,holonomy,
    spectrum,overlap,address,propagate,loops,conductance,potential,
    flux,wegscheider,check,network,species,reaction,solve,
    foreach,collect,rank,threshold,classify,scan,enumerate,
    synthesise,compare,select,where,report,batch},
  morekeywords=[2]{Coherence,DiseaseVector,Cell,Gap,Probe,Panel,
    Intervention,Dose,Reading,Verdict,Concern,Circuit,Holonomy,
    Spectrum,Overlap,Address,Loop,HolonomyVector,ROC,AUC,
    SpectrumDB,GrooveSpectrum,EscapeScore,Network},
  sensitive=true,
  morecomment=[l]{--},
  morestring=[b]",
  keywordstyle=\color{cfckw}\bfseries,
  keywordstyle=[2]\color{cfcnum},
  commentstyle=\color{cfccom}\itshape,
  stringstyle=\color{cfcstr},
  basicstyle=\ttfamily\small,
  backgroundcolor=\color{cfcbg},
  frame=single,framesep=4pt,xleftmargin=6pt,xrightmargin=4pt,
  breaklines=true,showstringspaces=false,columns=fullflexible,tabsize=2
}
\lstset{language=cfc}

% ---- Shorthand ----
\newcommand{\R}{\mathbb{R}}
\newcommand{\N}{\mathbb{N}}
\newcommand{\Rp}{\R_{\ge 0}}
\newcommand{\flo}{\beta}
\newcommand{\ecoh}{\eta}
\newcommand{\dvec}{\mathbf{D}}
\newcommand{\HolL}{H_\ell}
\newcommand{\Gcond}{\mathcal{G}}
\newcommand{\abs}[1]{\lvert#1\rvert}
\newcommand{\norm}[1]{\lVert#1\rVert}
\newcommand{\cf}[1]{\texttt{#1}}
\DeclareMathOperator*{\argmin}{arg\,min}
\DeclareMathOperator*{\argmax}{arg\,max}

\title{\textbf{Experimental Protocols for the Validation of\\
Circuit-Based Disease Diagnosis\\
via the Cause-for-Concern Calculus}\\[6pt]
\large Five Experiments on Public Data, with CFC Programs,\\
Syntax Extensions, and Falsification Criteria}

\author{%
Kundai Farai Sachikonye\\
\small Technical University of Munich\\
\small Chair of Proteomics and Bioanalytics\\
\small AIMe Registry for Artificial Intelligence\\
\small \texttt{kundai.sachikonye@bitspark.com}}

\date{\today}

\begin{document}
\maketitle

% =====================================================================
\begin{abstract}
\noindent
We design five experiments that test the principal claims of the
circuit-based disease framework against publicly available data, using the
Cause-for-Concern (CFC) domain-specific language as the implementation
vehicle. Each experiment is specified to the level at which it can be
executed without further design decisions: data source, extraction
protocol, CFC program, success criterion, failure criterion, and the
theorem it tests. Three syntax extensions to CFC are required and are
specified in full: a \cf{circuit} module for constructing and solving
metabolic circuits from reaction network databases, a \cf{spectrum}
module for Fourier-based spectral embeddings of biological sequences, and
a \cf{scan} form for batch evaluation over external databases. Each
extension is conservative: it introduces no new semantic primitive but
binds existing CFC operations to domain-specific data types. The five
experiments are: (E1)~holonomy-based diagnosis on the Recon3D human
metabolic reconstruction, testing whether loop inconsistency localises
enzyme defects without being told which enzyme is deficient;
(E2)~spectral MHC-peptide binding prediction on the Immune Epitope
Database, testing whether a zero-training-data spectral classifier
achieves competitive discrimination; (E3)~spectral codon-usage overlap
between human-infecting viruses and bacteriophages, testing the
housekeeping-mimicry prediction; (E4)~prospective immune escape
prediction on SARS-CoV-2, testing whether known Omicron mutations rank
in the top decile of spectral escape candidates; and (E5)~CFC reference
implementation on a five-node glycolytic circuit, validating the
ten-step protocol of the CFC specification. For each experiment we state
the quantitative threshold below which the result is negative and what
the negative result would falsify. No experiment requires wet-lab work,
proprietary data, or computational resources beyond a single workstation.
\end{abstract}

\tableofcontents

% =====================================================================
\section{Introduction}
\label{sec:intro}
% =====================================================================

\subsection{The validation gap}

The disease framework developed across the companion manuscripts rests on
a sequence of mathematical results: floor positivity, loop-holonomy
self-consistency, spectral matched-filter optimality, trajectory
certification, and gap-closure convergence. Every result is proved from
stated hypotheses, and every numerical check reported in those manuscripts
passes. The gap is that every check is either (a)~evaluated on synthetic
data generated under the same model the check assumes, or
(b)~consistent with known clinical orderings reproduced post hoc. No
experiment applies the framework to a dataset where the answer is not
already known and shows it outperforms an existing method or produces a
genuinely novel finding.

This paper closes the gap by designing five experiments that are
executable on public data, that test the framework's distinctive
predictions against established baselines, and that have quantitative
failure criteria stated in advance.

\subsection{Why CFC is the implementation language}

The Cause-for-Concern calculus \citep{sachikonye2026cfc} specifies a
language whose semantics is measurement, whose types enforce that no
therapy is applied without a localised cell, and whose operational
semantics is a monotone committed-measurement clock. The experiments
below are clinical-reasoning protocols---they localise a state, compute
a gap, and close it or decline---and CFC is the language designed to
express exactly that. Using CFC forces each experiment to satisfy the
accountability constraints the framework imposes: no zero-residue value,
no blind therapy, no unfounded allocation.

Three extensions are needed. They are specified in \Cref{sec:extensions}
and are conservative: each binds an existing CFC semantic operation
(probe, reading, localisation) to a domain-specific data type (circuit,
spectrum, database row) without introducing new semantic primitives.

\subsection{What this paper does and does not do}

It \emph{does} specify five experiments at sufficient detail to execute
them, including the CFC programs, the data extraction steps, and the
statistical tests. It \emph{does not} report results: no experiment has
been run. The paper is a protocol, not a report. It \emph{does} state
failure criteria: each experiment has a quantitative threshold below
which the result is negative and a statement of what the negative result
falsifies. It \emph{does not} claim that any experiment will succeed.

% =====================================================================
\section{CFC Syntax Extensions}
\label{sec:extensions}
% =====================================================================

The base CFC grammar of \citep{sachikonye2026cfc} exposes \cf{measure},
\cf{localize}, and \cf{close} over a physiological state system. The five
experiments below require three domain-specific modules that bind these
verbs to concrete data types. Each module is a conservative extension: it
adds keywords and type constructors but no new reduction rules. The
semantic model (\cf{measure} commits a reading, \cf{localize} shrinks a
surviving set, \cf{close} solves an $\ell_1$ program) is unchanged.

\subsection{The \texttt{circuit} module}
\label{sec:ext-circuit}

\begin{definition}[\cf{circuit} module: grammar extension]
\label{def:circuit-grammar}
The following productions are added to the CFC grammar.
\begin{align*}
\textit{circuitDecl} &\;::=\; \cf{circuit}\ \textit{ident}\
  \cf{from}\ \textit{expr}\ \cf{\{}\ \{\textit{circuitStmt}\}\ \cf{\}}\\[2pt]
\textit{circuitStmt} &\;::=\; \cf{species}\ \textit{ident}\
  \cf{:}\ \textit{speciesSpec}\\
  &\;\mid\; \cf{reaction}\ \textit{ident}\
  \cf{:}\ \textit{reactionSpec}\\
  &\;\mid\; \cf{solve}\ \cf{yield}\ \textit{ident}\\[2pt]
\textit{speciesSpec} &\;::=\; \cf{mu0}\ \cf{:}\ \textit{number}\
  \cf{,}\ \cf{concentration}\ \cf{:}\ \textit{number}\\[2pt]
\textit{reactionSpec} &\;::=\; \textit{ident}\ \cf{->}\ \textit{ident}\
  \cf{,}\ \cf{k}\ \cf{:}\ \textit{number}\\[2pt]
\textit{holonomyStmt} &\;::=\; \cf{holonomy}\ \cf{loops}\
  \cf{of}\ \textit{ident}\ \cf{yield}\ \textit{ident}\\[2pt]
\textit{wegscheiderStmt} &\;::=\; \cf{wegscheider}\ \cf{check}\
  \textit{ident}\ \cf{yield}\ \textit{ident}
\end{align*}
\end{definition}

\begin{definition}[\cf{circuit} module: types]
\label{def:circuit-types}
\leavevmode
\begin{itemize}[leftmargin=1.4em,itemsep=2pt]
\item \cf{Circuit}: a solved steady-state resistive circuit
  $\Gcond=(V,E,w)$ with node potentials $\mu_i$, conductances $G_{ij}$,
  and fluxes $J_{ij}$. Carries its floor annotation.
\item \cf{Loop}: a directed cycle in $\Gcond$, identified by its vertex
  sequence.
\item \cf{Holonomy}: the scalar $\norm{\HolL - \mathrm{Id}}_F$ for a
  single loop.
\item \cf{HolonomyVector}: the vector
  $\dvec_H = (\norm{\HolL - \mathrm{Id}}_F)_{\ell \in \mathcal{L}}$
  over all independent loops. This is a \cf{DiseaseVector} specialisation.
\end{itemize}
\end{definition}

\begin{definition}[\cf{circuit} module: semantics]
\label{def:circuit-sem}
\leavevmode
\begin{itemize}[leftmargin=1.4em,itemsep=2pt]
\item \cf{circuit C from data \{...\}} constructs $\Gcond$ by computing
  $\mu_i = \mu_i^\circ + RT \ln c_i$, $G_{ij} = k_{ij} c_i / (RT)$,
  $J_{ij} = G_{ij}(\mu_i - \mu_j)$ for each declared species and
  reaction. The \cf{solve} verb checks Kirchhoff's current law at every
  node and yields a \cf{Circuit} value. This is \cf{measure} specialised
  to the circuit domain: it commits the solved circuit as a reading.
\item \cf{holonomy loops of C yield H} enumerates all independent directed
  cycles in $\Gcond$, computes the Wegscheider sum
  $\sum_{\text{cycle}} \Delta\mu_{ij}$ for each, and yields the
  \cf{HolonomyVector}. This is \cf{localize} specialised to loop space:
  it shrinks the surviving set to the loops whose holonomy exceeds the
  floor.
\item \cf{wegscheider check C yield V} verifies
  $\abs{\sum_{\text{cycle}} \Delta\mu_{ij}} < \epsilon$ for every
  independent cycle, yielding a \cf{Verdict}. This is \cf{assert}
  specialised to thermodynamic consistency.
\end{itemize}
\end{definition}

\begin{remark}[Why these are conservative extensions]
\label{rem:conservative-circuit}
Each verb is a binding of an existing CFC semantic operation to a
domain-specific computation. \cf{solve} is \cf{measure} (it commits a
reading). \cf{holonomy loops} is \cf{localize} (it shrinks a surviving
set). \cf{wegscheider check} is \cf{assert} (it tests a condition and
may halt). The monotone committed-measurement clock advances at each
verb, the accountability type system applies, and no new reduction rule
is introduced. The extension is therefore sound by the soundness of the
base calculus (\citep{sachikonye2026cfc}, Theorem~10.1).
\end{remark}

\subsection{The \texttt{spectrum} module}
\label{sec:ext-spectrum}

\begin{definition}[\cf{spectrum} module: grammar extension]
\label{def:spectrum-grammar}
\begin{align*}
\textit{spectrumStmt} &\;::=\; \cf{synthesise}\ \cf{spectrum}\
  \cf{of}\ \textit{expr}\ [\cf{channels}\ \textit{number}]\
  [\cf{bins}\ \textit{number}]\ \cf{yield}\ \textit{ident}\\[2pt]
\textit{overlapStmt} &\;::=\; \cf{overlap}\ \textit{ident}\
  \textit{ident}\ \cf{yield}\ \textit{ident}\\[2pt]
\textit{batchStmt} &\;::=\; \cf{scan}\ \textit{ident}\
  \cf{against}\ \textit{expr}\\
  &\quad [\cf{where}\ \textit{cond}]\\
  &\quad \cf{collect}\ \textit{expr}\\
  &\quad \cf{yield}\ \textit{ident}
\end{align*}
\end{definition}

\begin{definition}[\cf{spectrum} module: types]
\label{def:spectrum-types}
\leavevmode
\begin{itemize}[leftmargin=1.4em,itemsep=2pt]
\item \cf{Spectrum}: an $\ell^2$-normalised vector
  $\psi(s) \in S^{cK-1}$ computed from a biological sequence $s$ by the
  channelised DFT embedding of \citep{sachikonye2026immuno}. Carries its
  floor annotation (the spectral Lipschitz constant bounds the
  resolution).
\item \cf{Overlap}: a scalar $\rho = \psi(q) \cdot \psi(t) \in [-1,1]$,
  the cosine similarity between two spectra. This is a \cf{Reading}
  specialisation.
\item \cf{GrooveSpectrum}: a \cf{Spectrum} computed from the 57
  polymorphic groove residues of an MHC allele. A \cf{Spectrum}
  subtype.
\item \cf{EscapeScore}: a scalar
  $E(v') = -\max_k \rho(\psi(v'), \psi(\mathrm{Ab}_k)) + w \cdot
  \rho(\psi(v'), \psi(R))$. A \cf{Reading} specialisation.
\end{itemize}
\end{definition}

\begin{definition}[\cf{spectrum} module: semantics]
\label{def:spectrum-sem}
\leavevmode
\begin{itemize}[leftmargin=1.4em,itemsep=2pt]
\item \cf{synthesise spectrum of seq channels c bins K yield psi}
  computes the channelised DFT: mean-subtract the $c \times L$ signal
  matrix, apply the real FFT to each channel, extract the first $K$
  magnitude bins, $\ell^2$-normalise. This is \cf{measure} specialised
  to spectral synthesis: it commits the spectrum as a reading. Cost:
  $\mathcal{O}(cL \log L)$.
\item \cf{overlap psi1 psi2 yield rho} computes the dot product
  $\psi_1 \cdot \psi_2$. This is a scalar \cf{Reading}. Cost:
  $\mathcal{O}(cK)$.
\item \cf{scan query against database where cond collect expr yield
  results} iterates over database entries, synthesises each spectrum on
  demand (empty database principle), evaluates \cf{cond}, and collects
  the expression. This is the batch form of \cf{measure}: each iteration
  commits one reading and advances the clock. The database is streamed;
  working memory is $\mathcal{O}(cK)$ per entry.
\end{itemize}
\end{definition}

\subsection{The \texttt{network} module}
\label{sec:ext-network}

\begin{definition}[\cf{network} module: grammar extension]
\label{def:network-grammar}
\begin{align*}
\textit{networkDecl} &\;::=\; \cf{network}\ \textit{ident}\
  \cf{from}\ \textit{expr}\\[2pt]
\textit{propagateStmt} &\;::=\; \cf{propagate}\ \textit{ident}\
  \cf{from}\ \textit{ident}\ \cf{yield}\ \textit{ident}\\[2pt]
\textit{addressStmt} &\;::=\; \cf{address}\ \textit{ident}\
  \cf{depth}\ \textit{number}\ \cf{yield}\ \textit{ident}
\end{align*}
\end{definition}

This module binds the CKG propagation operator $\mathcal{T}$ of
\citep{sachikonye2026ckg} to CFC's \cf{localize} semantics. A
\cf{propagate} call is a \cf{localize} restricted to the circuit's
contact graph: it returns the admissible set (a \cf{Cell} in CFC's
sense) of species causally reachable from a starting vertex under the
solved circuit's resting process. An \cf{address} call computes the
ternary interleaved address of a species from its solved-circuit
coordinate.

% =====================================================================
\section{Experiment 1: Holonomy Diagnosis on Recon3D}
\label{sec:exp1}
% =====================================================================

\subsection{Claim tested}

The Self-Consistency Primacy Theorem
(\citep{sachikonye2026blind}, Theorem~8.1): a cell is healthy if and only
if $\HolL = \mathrm{Id}$ for every reaction cycle $\ell$; disease is the
condition $\HolL \neq \mathrm{Id}$ for some $\ell$. The holonomy test
achieves perfect discrimination (AUC = 1.00 on synthetic circuits) and
\emph{localises} the defect to the broken loop without being told which
enzyme is deficient.

\subsection{Data source}

\begin{enumerate}[leftmargin=1.4em,itemsep=2pt]
\item \textbf{Recon3D} \citep{brunk2018recon3d}: the most comprehensive
  human metabolic reconstruction. We extract the glycolytic, pentose
  phosphate, and TCA cycle subnetwork ($\approx 80$ reactions,
  $\approx 60$ metabolites) with standard free energies $\Delta G^\circ$
  from the model's thermodynamic annotations.
\item \textbf{Healthy erythrocyte fluxes}: Rapoport \& Heinrich (1975),
  Mulquiney \& Kuchel (1999). Measured glycolytic flux
  $J_{\mathrm{glycolysis}} = 1.1 \pm 0.1$ mmol/L/h; metabolite
  concentrations from physiological reference ranges.
\item \textbf{Pyruvate kinase (PK) deficiency}: Zanella et al.\ (2005).
  PK activity reduced to $5$--$30\%$ of normal; upstream metabolite
  accumulation (2,3-BPG, PEP) and downstream depletion (pyruvate, ATP)
  documented.
\item \textbf{Cancer metabolic profiles}: Jain et al.\ (2012), NCI-60
  panel. Metabolic consumption and release rates for 60 cancer cell
  lines across 219 metabolites.
\end{enumerate}

\subsection{Protocol}

\begin{protocol}[Holonomy diagnosis]
\label{prot:holonomy}
\leavevmode
\begin{enumerate}[leftmargin=1.4em,itemsep=3pt]
\item \textbf{Circuit construction.} For each dataset (healthy, PK-deficient,
  cancer), construct the circuit graph $\Gcond = (V, E, w)$:
  \begin{itemize}[leftmargin=1.2em,itemsep=1pt]
  \item Node potential: $\mu_i = \mu_i^\circ + RT \ln c_i$ with
    $RT = 2.577$ kJ/mol at $T = 310$ K.
  \item Edge conductance: $G_{ij} = k_{ij} c_i / (RT)$.
  \item Flux: $J_{ij} = G_{ij}(\mu_i - \mu_j)$.
  \end{itemize}
\item \textbf{KCL verification.} At every internal node, verify
  $\abs{\sum J_{\mathrm{in}} - \sum J_{\mathrm{out}}} < \epsilon$ with
  $\epsilon = 10^{-6}$. Any node failing this check invalidates the
  circuit and halts the experiment.
\item \textbf{Cycle enumeration.} Enumerate all independent directed
  cycles in the subnetwork using the cycle space basis of the directed
  graph (rank of the cycle matrix = $\abs{E} - \abs{V} + 1$).
\item \textbf{Holonomy computation.} For each cycle $\ell$, compute the
  Wegscheider sum:
  \begin{equation}
  \HolL = \sum_{(i,j) \in \ell} (\mu_j - \mu_i).
  \label{eq:holonomy-sum}
  \end{equation}
  By the voltage law (\citep{sachikonye2026ckg}, Theorem~2.1), this
  telescopes to zero for a consistent circuit. Disease is
  $\abs{\HolL} > \flo$ for some $\ell$.
\item \textbf{Classifier construction.} Define two classifiers:
  \begin{itemize}[leftmargin=1.2em,itemsep=1pt]
  \item \emph{Holonomy classifier}: threshold on
    $\max_\ell \abs{\HolL}$.
  \item \emph{Template classifier}: Euclidean distance
    $\norm{\mathbf{c}_{\mathrm{patient}} -
    \mathbf{c}_{\mathrm{healthy}}}$ from the healthy metabolite
    concentration vector.
  \end{itemize}
\item \textbf{Discrimination.} Compute ROC-AUC for both classifiers on
  the healthy-vs-diseased discrimination task.
\item \textbf{Localisation.} For PK-deficient circuits, identify which
  loop(s) have $\abs{\HolL} > \flo$. The prediction is that the ATP
  production loop (containing the PK step) is flagged, and the pentose
  phosphate pathway loops are not.
\end{enumerate}
\end{protocol}

\subsection{CFC program}

\begin{lstlisting}[caption={E1: holonomy diagnosis on Recon3D.}]
-- exp1_holonomy.cfc
floor 1e-6

import Recon3D              -- Recon3D thermodynamic annotations
import Erythrocyte.Healthy   -- Rapoport & Heinrich concentrations
import Erythrocyte.PKDeficient -- Zanella et al. concentrations

-- construct the healthy circuit
circuit Healthy from Recon3D.glycolysis_tca {
  solve yield C_healthy
}

-- construct the PK-deficient circuit
circuit PKDef from Recon3D.glycolysis_tca {
  -- override PK activity to 15% of normal
  reaction PK : PEP -> Pyruvate, k : 0.15 * Recon3D.k_PK
  solve yield C_pkdef
}

-- verify Kirchhoff current law on both
wegscheider check C_healthy yield V_h
assert V_h == resolved emit "healthy circuit consistent"

wegscheider check C_pkdef yield V_pk
-- PK-deficient circuit may or may not pass Wegscheider;
-- failure is itself a finding

-- compute holonomy on all independent loops
holonomy loops of C_healthy yield H_healthy
holonomy loops of C_pkdef yield H_pkdef

-- the diagnostic question: which loops are inconsistent?
patient Erythrocyte_PKDef {
  -- localise using the holonomy vector as the probe reading
  measure panel { holonomy loops of C_pkdef } yield reading

  localize reading
    using panel { holonomy loops of C_pkdef }
    until resolved otherwise decline
    yield here

  -- here.disease_vector is the holonomy vector;
  -- its nonzero entries are the broken loops
  let broken_loops := here.disease_vector
    where abs(H) > floor

  -- the prediction: PK-containing loop is broken
  assert broken_loops contains "ATP_production_loop"
    emit "localisation succeeded: PK loop identified"

  -- gap: displacement to the nearest consistent circuit
  close gap of here
    toward coherent
    using basis Recon3D.enzymes
    until coherent otherwise decline
    yield plan

  -- plan.intervention names the enzyme(s) whose conductance
  -- modification restores loop consistency
  report plan
}
\end{lstlisting}

\subsection{Success and failure criteria}

\begin{criterion}[E1 success]
\label{crit:e1-success}
\leavevmode
\begin{enumerate}[label=(\alph*),leftmargin=1.8em,itemsep=2pt]
\item The healthy circuit has $\abs{\HolL} < 10^{-12}$ on every
  independent cycle (machine-precision zero).
\item The PK-deficient circuit has $\abs{\HolL} > 10^{-3}$ on at
  least one cycle containing the PK reaction step.
\item The holonomy classifier achieves ROC-AUC $\geq 0.95$ on the
  healthy-vs-PK-deficient discrimination.
\item The holonomy vector localises the defect to the PK-containing
  loop(s) without being told which enzyme is deficient.
\end{enumerate}
\end{criterion}

\begin{criterion}[E1 failure]
\label{crit:e1-failure}
If the healthy circuit has $\abs{\HolL} > 10^{-6}$ on any cycle, the
circuit construction is incorrect and the experiment is invalid (not
negative). If the PK-deficient circuit has $\abs{\HolL} < 10^{-6}$ on
all cycles, the holonomy test does not detect the defect, and the
Self-Consistency Primacy Theorem is empirically unsupported on real
metabolic data. If the holonomy classifier achieves AUC $< 0.80$, the
method does not outperform template comparison, and the claimed
superiority of self-consistency over template-based diagnosis is not
confirmed.
\end{criterion}

\begin{remark}[What a negative result would mean]
A negative result on criterion (b) would not falsify the theorem, which
is a mathematical identity about Kirchhoff's voltage law. It would show
that the concentrations and rate constants available from public databases
are not accurate enough to produce a detectable holonomy defect, which is
a limitation of the data, not of the theory. The experiment is designed
to distinguish these two cases: criterion (a) tests whether the healthy
circuit is thermodynamically consistent (a property of the data quality),
and only if (a) passes does (b) test the theory.
\end{remark}

% =====================================================================
\section{Experiment 2: Spectral MHC Binding on IEDB}
\label{sec:exp2}
% =====================================================================

\subsection{Claim tested}

The molecular recognition framework
(\citep{sachikonye2026immuno}, Definition~3.2): a peptide $p$ binds MHC
allele $M$ if and only if
$\rho(p, M) \equiv \psi(p) \cdot \Gamma(M) > \rho^*_{\mathrm{MHC}}$.
The spectral matched filter is Neyman-Pearson optimal
(\citep{sachikonye2026immuno}, Theorem~2.1) and requires zero training
data.

\subsection{Data source}

The Immune Epitope Database (IEDB; \texttt{iedb.org}):
\begin{itemize}[leftmargin=1.4em,itemsep=2pt]
\item Download: all human HLA class I peptide-MHC binding assays with
  IC$_{50}$ measurements.
\item Filter: 9-mer peptides only; HLA-A, -B, -C alleles with $\geq 200$
  measurements each.
\item Expected yield: $\approx 500{,}000$ peptide-allele pairs across
  $\approx 100$ alleles.
\item Binarisation: IC$_{50} \leq 500$ nM $\Rightarrow$ binder (positive);
  IC$_{50} > 500$ nM $\Rightarrow$ non-binder (negative). This is the
  standard threshold used by NetMHCpan \citep{reynisson2020}.
\end{itemize}

HLA groove residue sequences: IPD-IMGT/HLA database
(\texttt{ebi.ac.uk/ipd/imgt/hla/}). For each allele $M$, extract the 57
polymorphic binding-groove residues of the $\alpha_1$ and $\alpha_2$
domains.

\subsection{Protocol}

\begin{protocol}[Spectral MHC binding]
\label{prot:mhc}
\leavevmode
\begin{enumerate}[leftmargin=1.4em,itemsep=3pt]
\item \textbf{Groove spectrum computation.} For each HLA allele $M$:
  \begin{enumerate}[label=(\roman*),leftmargin=1.8em]
  \item Extract the 57 groove residues from IPD-IMGT/HLA.
  \item Compute the three-channel protein signal matrix
    $X \in \R^{3 \times 57}$: hydropathy (Kyte-Doolittle, normalised to
    $[0,1]$), van der Waals volume (normalised to $[0,1]$), net charge
    ($\{-1, 0, +1\}$).
  \item Mean-subtract each channel.
  \item Apply the real FFT; retain the first $K = 16$ non-DC magnitude
    bins per channel.
  \item $\ell^2$-normalise to obtain
    $\Gamma(M) \in S^{47}$.
  \end{enumerate}
\item \textbf{Peptide spectrum computation.} For each 9-mer peptide $p$,
  compute $\psi(p) \in S^{47}$ identically.
\item \textbf{Binding score.} For each peptide-allele pair, compute
  $\rho = \psi(p) \cdot \Gamma(M)$.
\item \textbf{Per-allele ROC-AUC.} For each allele with $\geq 200$
  measurements, compute the ROC-AUC of $\rho$ as a classifier against
  the IC$_{50} \leq 500$ nM binarisation.
\item \textbf{Comparison.} Compare per-allele AUC against published
  NetMHCpan 4.1 benchmark AUCs \citep{reynisson2020}.
\item \textbf{Rare-allele analysis.} Stratify alleles by IEDB
  measurement count. Compute the mean AUC in each stratum for both the
  spectral method and NetMHCpan. The prediction is that the spectral
  method's AUC is allele-frequency-independent (because it uses zero
  training data), while NetMHCpan degrades on rare alleles.
\item \textbf{Supertype validation.} Compute
  $\Gamma(M_i) \cdot \Gamma(M_j)$ for all allele pairs. Cluster at
  cosine similarity $\geq 0.85$. Compare clusters against the Sette \&
  Sidney (1999) HLA supertypes.
\end{enumerate}
\end{protocol}

\subsection{CFC program}

\begin{lstlisting}[caption={E2: spectral MHC binding on IEDB.}]
-- exp2_mhc_binding.cfc
floor 1e-3

import IEDB                  -- peptide-MHC binding assays
import IPD_IMGT_HLA          -- HLA allele groove sequences

-- compute groove spectra for all alleles
basis HLA_Grooves {
  probe groove_spectrum(allele) {
    reference: IPD_IMGT_HLA.groove_residues(allele)
    observe: synthesise spectrum of reference
               channels 3 bins 16
  }
}

-- the experiment: scan all IEDB peptide-allele pairs
patient IEDB_Benchmark {
  -- for each allele with >= 200 measurements
  foreach allele in IPD_IMGT_HLA.alleles
    where IEDB.count(allele) >= 200 {

    -- compute groove spectrum
    synthesise spectrum of
      IPD_IMGT_HLA.groove_residues(allele)
      channels 3 bins 16
      yield groove

    -- scan all peptides for this allele
    scan groove against IEDB.peptides(allele)
      where IEDB.peptide_length == 9
      collect overlap(groove, synthesise spectrum of
                IEDB.peptide channels 3 bins 16)
      yield scores

    -- compute ROC-AUC against IC50 <= 500nM
    let auc := ROC(scores, IEDB.binders(allele, 500))
    report allele, auc
  }

  -- rare-allele stratification
  let rare_alleles := IPD_IMGT_HLA.alleles
    where IEDB.count(allele) < 50
  let common_alleles := IPD_IMGT_HLA.alleles
    where IEDB.count(allele) >= 1000

  report "rare_mean_auc", mean(auc for rare_alleles)
  report "common_mean_auc", mean(auc for common_alleles)
}
\end{lstlisting}

\subsection{Success and failure criteria}

\begin{criterion}[E2 success]
\label{crit:e2-success}
\leavevmode
\begin{enumerate}[label=(\alph*),leftmargin=1.8em,itemsep=2pt]
\item Mean per-allele AUC $\geq 0.80$ across all alleles with
  $\geq 200$ measurements.
\item Mean AUC on rare alleles ($< 50$ IEDB measurements) is within
  $0.05$ of mean AUC on common alleles ($\geq 1000$ measurements).
\item HLA supertype clusters recovered at cosine $\geq 0.85$ match the
  Sette \& Sidney supertypes with adjusted Rand index $\geq 0.60$.
\end{enumerate}
\end{criterion}

\begin{criterion}[E2 failure]
\label{crit:e2-failure}
If mean AUC $< 0.70$, the magnitude-only spectral embedding is too
lossy for MHC binding prediction, and the phase-discarding design of the
spectral fingerprint is falsified as a binding predictor. If the rare/common
AUC gap exceeds $0.10$ in favour of common alleles, the spectral method
is not allele-frequency-independent and the zero-training-data advantage
is not confirmed.
\end{criterion}

% =====================================================================
\section{Experiment 3: Spectral Codon-Usage Overlap}
\label{sec:exp3}
% =====================================================================

\subsection{Claim tested}

The housekeeping-mimicry prediction
(\citep{sachikonye2026viral}, Proposition~5.1): viral RNA genomes must
present charge-oscillation profiles matching the host's constitutively
expressed housekeeping genes, because those are the templates the host
machinery is calibrated to process. Cross-kingdom non-overlap
(\citep{sachikonye2026viral}, Proposition~8.1): bacteriophage proteins
have zero spectral overlap with human housekeeping centroids because
phage codons are tuned to bacterial ribosomal decoding.

\subsection{Data source}

\begin{enumerate}[leftmargin=1.4em,itemsep=2pt]
\item \textbf{Human housekeeping genes.} Eisenberg \& Levanon (2013)
  consensus list: $\approx 3{,}800$ genes. CDS sequences from Ensembl
  GRCh38.
\item \textbf{Human-infecting viral proteins.} UniProt reviewed entries
  for HIV-1, SARS-CoV-2, Influenza A (H1N1, H3N2), HCV, EBV, HSV-1,
  HPV-16: $\approx 200$ proteins total.
\item \textbf{Bacteriophage proteins.} PhagesDB and NCBI phage genomes:
  T4, $\lambda$, P22, T7, $\phi$X174, $\approx 300$ proteins.
\item \textbf{E.\ coli housekeeping genes.} Baba et al.\ (2006), Keio
  collection essential genes: $\approx 300$ genes.
\end{enumerate}

\subsection{Protocol}

\begin{protocol}[Codon-usage spectral overlap]
\label{prot:codon}
\leavevmode
\begin{enumerate}[leftmargin=1.4em,itemsep=3pt]
\item \textbf{Compute housekeeping centroids.}
  \begin{enumerate}[label=(\roman*),leftmargin=1.8em]
  \item For each human housekeeping protein, compute the 3-channel
    spectral embedding $\psi(g_k)$ with $K = 16$.
  \item Compute
    $\Gamma_{\mathrm{HK}}^{\mathrm{human}} =
    \frac{1}{N}\sum_k \psi(g_k) /
    \norm{\frac{1}{N}\sum_k \psi(g_k)}$.
  \item Repeat for E.\ coli housekeeping genes to obtain
    $\Gamma_{\mathrm{HK}}^{E.\,coli}$.
  \end{enumerate}
\item \textbf{Compute viral and phage spectra.} For each protein $v_j$
  (human virus) or $p_j$ (phage), compute $\psi(v_j)$ or $\psi(p_j)$.
\item \textbf{Compute overlap scores.}
  \begin{itemize}[leftmargin=1.2em,itemsep=1pt]
  \item $\rho_{\mathrm{human}}(v_j) =
    \psi(v_j) \cdot \Gamma_{\mathrm{HK}}^{\mathrm{human}}$
    for each human virus protein.
  \item $\rho_{\mathrm{human}}(p_j) =
    \psi(p_j) \cdot \Gamma_{\mathrm{HK}}^{\mathrm{human}}$
    for each phage protein.
  \item $\rho_{E.\,coli}(v_j) =
    \psi(v_j) \cdot \Gamma_{\mathrm{HK}}^{E.\,coli}$
    for each human virus protein.
  \item $\rho_{E.\,coli}(p_j) =
    \psi(p_j) \cdot \Gamma_{\mathrm{HK}}^{E.\,coli}$
    for each phage protein.
  \end{itemize}
\item \textbf{Statistical test.} Wilcoxon rank-sum test comparing
  $\{\rho_{\mathrm{human}}(v_j)\}$ against
  $\{\rho_{\mathrm{human}}(p_j)\}$: the prediction is that human virus
  proteins have significantly higher overlap with human housekeeping
  centroids than phage proteins do ($p < 0.001$). Symmetrically, phage
  proteins should have higher overlap with E.\ coli centroids.
\item \textbf{Visualisation.} Plot each protein in the
  $(\rho_{\mathrm{human}}, \rho_{E.\,coli})$ plane. The prediction is
  that human viruses cluster in the high-$\rho_{\mathrm{human}}$ region
  and phages in the high-$\rho_{E.\,coli}$ region, with clean
  separation.
\end{enumerate}
\end{protocol}

\subsection{CFC program}

\begin{lstlisting}[caption={E3: spectral codon-usage overlap.}]
-- exp3_codon_overlap.cfc
floor 1e-3

import Ensembl.GRCh38        -- human CDS sequences
import Eisenberg2013         -- housekeeping gene list
import UniProt.Viral         -- human virus proteins
import PhagesDB              -- bacteriophage proteins
import Baba2006              -- E. coli essential genes

-- compute human housekeeping centroid
let hk_human := Eisenberg2013.genes
scan each in hk_human
  collect synthesise spectrum of
    Ensembl.CDS(each) channels 3 bins 16
  yield hk_spectra_human

let centroid_human := normalise(mean(hk_spectra_human))

-- compute E. coli housekeeping centroid
let hk_ecoli := Baba2006.essential_genes
scan each in hk_ecoli
  collect synthesise spectrum of each channels 3 bins 16
  yield hk_spectra_ecoli

let centroid_ecoli := normalise(mean(hk_spectra_ecoli))

-- scan human virus proteins
scan v in UniProt.Viral.human_infecting
  collect {
    rho_human: overlap(synthesise spectrum of v
                 channels 3 bins 16, centroid_human),
    rho_ecoli: overlap(synthesise spectrum of v
                 channels 3 bins 16, centroid_ecoli)
  }
  yield virus_scores

-- scan phage proteins
scan p in PhagesDB.proteins
  collect {
    rho_human: overlap(synthesise spectrum of p
                 channels 3 bins 16, centroid_human),
    rho_ecoli: overlap(synthesise spectrum of p
                 channels 3 bins 16, centroid_ecoli)
  }
  yield phage_scores

-- statistical test
assert wilcoxon(virus_scores.rho_human,
                phage_scores.rho_human).p < 0.001
  emit "human viruses have higher human-HK overlap"

assert wilcoxon(phage_scores.rho_ecoli,
                virus_scores.rho_ecoli).p < 0.001
  emit "phages have higher ecoli-HK overlap"

report virus_scores, phage_scores
\end{lstlisting}

\subsection{Success and failure criteria}

\begin{criterion}[E3 success]
\label{crit:e3-success}
\leavevmode
\begin{enumerate}[label=(\alph*),leftmargin=1.8em,itemsep=2pt]
\item Mean $\rho_{\mathrm{human}}$ for human virus proteins exceeds mean
  $\rho_{\mathrm{human}}$ for phage proteins by $> 0.05$, with Wilcoxon
  $p < 0.001$.
\item Mean $\rho_{E.\,coli}$ for phage proteins exceeds mean
  $\rho_{E.\,coli}$ for human virus proteins by $> 0.05$, with Wilcoxon
  $p < 0.001$.
\item In the $(\rho_{\mathrm{human}}, \rho_{E.\,coli})$ plane, the two
  populations are separable by a linear classifier with AUC $\geq 0.80$.
\end{enumerate}
\end{criterion}

\begin{criterion}[E3 failure]
If the distributions overlap substantially (Wilcoxon $p > 0.05$ on either
test), the spectral embedding does not capture the host-specificity
signature that the housekeeping-mimicry hypothesis predicts, and the
cross-kingdom non-overlap claim is empirically unsupported.
\end{criterion}

% =====================================================================
\section{Experiment 4: Prospective Escape Prediction}
\label{sec:exp4}
% =====================================================================

\subsection{Claim tested}

The immune escape landscape
(\citep{sachikonye2026immuno}, Definition~5.1;
\citep{sachikonye2026pathogen}, Algorithm~2): escape variants are
mutations that shift the pathogen spectrum out of antibody recognition
bands while retaining receptor-binding overlap. Known Omicron mutations
should rank in the top decile of predicted escape candidates.

\subsection{Data source}

\begin{enumerate}[leftmargin=1.4em,itemsep=2pt]
\item \textbf{Spike RBD sequence.} Wuhan-Hu-1 (GenBank MN908947),
  residues 319--541 (223 amino acids).
\item \textbf{ACE2 receptor.} Human ACE2 ectodomain (UniProt Q9BYF1),
  residues 19--615.
\item \textbf{Pre-Omicron antibody panel.} Cameroni et al.\ (2022,
  PMID 35016198): heavy and light chain CDR sequences for S309, S2E12,
  S2H58, S2K146, S2X259, S2D106 ($\geq 6$ antibodies).
\item \textbf{Ground truth.} Omicron BA.1 RBD mutations: K417N, N440K,
  G446S, S477N, T478K, E484A, Q493R, G496S, Q498R, N501Y, Y505H (11
  substitutions).
\end{enumerate}

\subsection{Protocol}

\begin{protocol}[Escape prediction]
\label{prot:escape}
\leavevmode
\begin{enumerate}[leftmargin=1.4em,itemsep=3pt]
\item \textbf{Receptor and antibody spectra.} Compute $\psi(\mathrm{ACE2})$
  and $\Gamma(\mathrm{Ab}_k)$ for each antibody (from concatenated CDR
  sequences) using 3-channel protein embedding, $K = 16$.
\item \textbf{Variant enumeration.} Enumerate all $19 \times 223 = 4{,}237$
  single-amino-acid substitutions in the RBD.
\item \textbf{Escape scoring.} For each variant $v'$:
  \begin{enumerate}[label=(\roman*),leftmargin=1.8em]
  \item Synthesise $\psi(v')$.
  \item $\rho_{\mathrm{bind}} = \psi(v') \cdot \psi(\mathrm{ACE2})$.
  \item $\rho_{\mathrm{escape}} = \max_k \psi(v') \cdot \Gamma(\mathrm{Ab}_k)$.
  \item $E(v') = -\rho_{\mathrm{escape}} + 0.8 \cdot \rho_{\mathrm{bind}}$.
  \end{enumerate}
\item \textbf{Ranking.} Sort all 4{,}237 variants by $E(v')$ descending.
\item \textbf{Evaluation.} For each of the 11 Omicron BA.1 mutations,
  record its percentile rank in the sorted list.
\item \textbf{Statistical test.} Under the null hypothesis that Omicron
  mutations are uniformly distributed, compute the probability that all
  11 fall in the top $p$-percentile: $P = p^{11}$. Report the minimum
  $p$ at which the null is rejected at $\alpha = 0.01$.
\item \textbf{Control.} Repeat with a random ``antibody'' spectrum
  (uniform on $S^{47}$) replacing the real panel. The enrichment should
  vanish.
\end{enumerate}
\end{protocol}

\subsection{CFC program}

\begin{lstlisting}[caption={E4: SARS-CoV-2 escape prediction.}]
-- exp4_escape.cfc
floor 1e-3

import GenBank.MN908947       -- Wuhan-Hu-1 Spike
import UniProt.Q9BYF1         -- human ACE2
import Cameroni2022           -- antibody CDR sequences

let RBD := GenBank.MN908947.spike[319:541]

-- compute receptor spectrum
synthesise spectrum of UniProt.Q9BYF1.ectodomain
  channels 3 bins 16 yield psi_ACE2

-- compute antibody panel spectra
foreach ab in Cameroni2022.antibodies {
  synthesise spectrum of ab.CDR_concatenated
    channels 3 bins 16 yield psi_ab
}

-- enumerate all single-AA substitutions
let variants := enumerate_substitutions(RBD)
  -- yields 19 * 223 = 4237 variants

-- score each variant
scan v in variants
  collect {
    mutation: v.label,
    psi_v: synthesise spectrum of v.sequence
             channels 3 bins 16,
    rho_bind: overlap(psi_v, psi_ACE2),
    rho_escape: max(
      foreach ab in Cameroni2022.antibodies
        collect overlap(psi_v, psi_ab)),
    E: -rho_escape + 0.8 * rho_bind
  }
  yield escape_landscape

-- rank by escape score
let ranked := rank(escape_landscape, by: E, descending)

-- evaluate Omicron BA.1 mutations
let omicron_mutations := [
  "K417N","N440K","G446S","S477N","T478K",
  "E484A","Q493R","G496S","Q498R","N501Y","Y505H"
]

foreach m in omicron_mutations {
  let percentile := ranked.percentile(m)
  report m, percentile
}

-- statistical test
let min_percentile := max(
  foreach m in omicron_mutations collect ranked.percentile(m))
assert min_percentile <= 0.10
  emit "all Omicron mutations in top 10%"

-- negative control: random antibody panel
let random_panel := uniform_sphere(48, seed: 42)
scan v in variants
  collect {
    E_random: -overlap(psi_v, random_panel) + 0.8 * rho_bind
  }
  yield random_landscape

report "real_enrichment", enrichment(ranked, omicron_mutations)
report "random_enrichment",
  enrichment(rank(random_landscape), omicron_mutations)
\end{lstlisting}

\subsection{Success and failure criteria}

\begin{criterion}[E4 success]
\label{crit:e4-success}
\leavevmode
\begin{enumerate}[label=(\alph*),leftmargin=1.8em,itemsep=2pt]
\item At least $8/11$ Omicron BA.1 mutations rank in the top $10\%$ of
  escape candidates.
\item The median percentile of Omicron mutations is $\leq 15\%$.
\item The enrichment under the real antibody panel is at least $3\times$
  higher than under the random control.
\end{enumerate}
\end{criterion}

\begin{criterion}[E4 failure]
If fewer than $6/11$ Omicron mutations fall in the top $20\%$, the
spectral escape landscape does not capture real immunological selection
pressure, and the claim that escape variants are spectrally predictable
is empirically unsupported. If the random control achieves comparable
enrichment, the result is an artefact of the scoring function geometry
rather than a biological signal.
\end{criterion}

% =====================================================================
\section{Experiment 5: CFC Reference Implementation}
\label{sec:exp5}
% =====================================================================

\subsection{Claim tested}

The CFC specification (\citep{sachikonye2026cfc}): the ten-step
validation protocol of \S13, establishing that the Floor Theorem,
Localization Theorem, Certificate Theorem, Gap-Closure Theorem, and
Wellness Theorem hold on an explicit model instance.

\subsection{Model}

A five-node glycolytic circuit:
\begin{equation}
\text{Glucose} \xrightarrow{k_1} \text{G6P}
\xrightarrow{k_2} \text{FBP}
\xrightarrow{k_3} \text{G3P}
\xrightarrow{k_4} \text{Pyruvate}
\label{eq:glycolysis-model}
\end{equation}
with one feedback loop: ATP/ADP ratio coupled to the hexokinase ($k_1$)
and phosphofructokinase ($k_2$) steps. Mass-action kinetics with
textbook rate constants from Teusink et al.\ (2000).

Three oscillator classes: glycolytic flux ($\Pi_1 = J_{\mathrm{total}}$),
ATP/ADP ratio ($\Pi_2 = [\mathrm{ATP}]/[\mathrm{ADP}]$), redox state
($\Pi_3 = [\mathrm{NADH}]/[\mathrm{NAD}^+]$).

\subsection{Protocol}

\begin{protocol}[CFC reference implementation]
\label{prot:cfc}
The ten validation steps from \citep{sachikonye2026cfc}, \S13:
\begin{enumerate}[leftmargin=1.4em,itemsep=3pt]
\item \textbf{Floor.} Set readout resolution $\delta = 10^{-3}$.
  Confirm the induced state-discrimination has $\flo > 0$. Perturb
  states by $\Delta x < \flo$ and confirm readouts are
  indistinguishable. Perturb by $\Delta x > \flo$ and confirm readouts
  differ.
\item \textbf{Coherence.} Compute $\ecoh_i$ for each oscillator class
  at optimal and degraded values. Confirm $\ecoh_i \in [0,1]$. Confirm
  the coherent set $\mathcal{M}^*_\theta$ equals the states with all
  $\ecoh_i \geq \theta$.
\item \textbf{Localization.} Construct a glucose-challenge probe
  (ingestion of 75 units glucose, observe G6P at $t = 10$). Confirm it
  is floor-injective on the glycolytic subspace: distinct metabolic
  states produce distinct traces. Construct a probe that is \emph{not}
  floor-injective (observe total carbon only) and confirm it fails to
  localise.
\item \textbf{Certificate.} Build a cascade probe:
  glucose $\to$ (state-gated glycolysis) $\to$ pyruvate $\to$
  (state-gated LDH) $\to$ lactate. Confirm lactate presence certifies
  the entire upstream path. Block LDH ($k_{\mathrm{LDH}} = 0$); confirm
  lactate is absent.
\item \textbf{Rigidity.} Grow the observed species set $S$ and confirm
  $\mathcal{M}_S$ is non-increasing. Confirm it becomes cell-tight
  ($\mathrm{diam}(\mathcal{M}_S) < \flo$) after a panel of $3$ probes.
  Confirm that with only $1$ probe, $\mathcal{M}_S$ is not cell-tight
  and \cf{localize} returns \cf{decline}.
\item \textbf{Gap-closure.} Place the model in a PK-deficient state
  ($k_4 = 0.15 k_4^{\mathrm{normal}}$). Compute the gap
  $\mathbf{g} = x^* - x$ to the nearest coherent state. Solve the
  $\ell_1$ program over the enzyme basis. Confirm the selected
  intervention is sparse ($\leq 3$ nonzero doses). Iterate the
  contraction; confirm $\mathrm{dist}(x_k, \mathcal{M}^*) \leq
  \lambda^k \mathrm{dist}(x_0, \mathcal{M}^*)$ with $\lambda < 1$.
\item \textbf{Unity.} Confirm the glucose probe simultaneously
  localises (its trace pins the state to a cell) and displaces (its
  metabolic products shift the state).
\item \textbf{Decline.} Construct a model variant with two metabolically
  distinct states (aerobic vs.\ anaerobic) that produce identical
  glucose-challenge readouts. Confirm \cf{localize} returns
  \cf{contested decline} and names the two surviving cells.
\item \textbf{Derived vs.\ fitted.} Fit a polynomial surrogate on
  $[\mathrm{ATP}] \in [0.5, 2.0]$ mM. Evaluate both the derived
  (mass-action) and surrogate models at $[\mathrm{ATP}] = 5.0$ mM
  (off-training). Confirm the surrogate error exceeds the derived-model
  error by $> 10\times$.
\item \textbf{Type safety.} Confirm the type checker rejects: (a)~a
  literal of zero floor; (b)~a \cf{close} whose subject is a
  \cf{Reading}, not a \cf{Cell}; (c)~a probe with floor annotation
  $\flo = 0$.
\end{enumerate}
\end{protocol}

\subsection{CFC program}

\begin{lstlisting}[caption={E5: CFC reference implementation.}]
-- exp5_reference.cfc
floor 1e-3

-- the five-node glycolytic model
circuit Glycolysis from Teusink2000 {
  species Glucose     : mu0 : -917.0, concentration : 5.0
  species G6P         : mu0 : -1760.0, concentration : 0.5
  species FBP         : mu0 : -2600.0, concentration : 0.1
  species G3P         : mu0 : -1510.0, concentration : 0.05
  species Pyruvate    : mu0 : -474.0, concentration : 0.1

  reaction HK  : Glucose -> G6P,    k : 0.1
  reaction PFK : G6P -> FBP,        k : 0.05
  reaction Ald : FBP -> G3P,        k : 0.08
  reaction PK  : G3P -> Pyruvate,   k : 0.12

  solve yield C_normal
}

-- V1: floor positivity
patient V1_Floor {
  measure panel { glucose_challenge(75) } yield r1
  let perturbed := perturb(C_normal, G6P, delta: 0.5e-3)
  measure panel { glucose_challenge(75) } yield r2
  assert abs(r1 - r2) < floor
    emit "V1: sub-floor perturbation indistinguishable"

  let perturbed2 := perturb(C_normal, G6P, delta: 5e-3)
  measure panel { glucose_challenge(75) } yield r3
  assert abs(r1 - r3) > floor
    emit "V1: supra-floor perturbation distinguishable"
}

-- V2: coherence
patient V2_Coherence {
  let eta_flux := coherence(C_normal, "glycolytic_flux",
    optimal: 1.1, degraded: 0.0)
  assert eta_flux >= 0.0 emit "V2: coherence >= 0"
  assert eta_flux <= 1.0 emit "V2: coherence <= 1"
}

-- V3: localization
patient V3_Localize {
  measure panel { glucose_challenge(75) >> observe(G6P) }
    yield r_injective

  localize r_injective
    using panel { glucose_challenge(75) >> observe(G6P) }
    until resolved otherwise decline
    yield cell_or_decline

  assert cell_or_decline == resolved
    emit "V3: floor-injective probe localises"
}

-- V5: rigidity
patient V5_Rigidity {
  -- single probe: should decline
  localize baseline
    using panel { glucose_challenge(75) }
    until resolved otherwise decline
    yield result_1probe

  assert result_1probe == decline
    emit "V5: single probe declines (correct)"

  -- three probes: should resolve
  localize baseline
    using panel {
      glucose_challenge(75) >> observe(G6P)
      || lactate_probe(10)
      || atp_ratio_probe()
    }
    until resolved otherwise decline
    yield result_3probes

  assert result_3probes == resolved
    emit "V5: three probes resolve"
}

-- V6: gap-closure
patient V6_GapClosure {
  -- PK-deficient state
  circuit PKDef from Glycolysis {
    reaction PK : G3P -> Pyruvate, k : 0.15 * 0.12
    solve yield C_pkdef
  }

  localize C_pkdef
    using panel {
      glucose_challenge(75) >> observe(G6P)
      || lactate_probe(10)
      || atp_ratio_probe()
    }
    until resolved
    yield here

  close gap of here
    toward coherent
    using basis Glycolysis.enzymes
    until coherent otherwise decline
    yield plan

  -- confirm sparsity
  assert plan.nonzero_doses <= 3
    emit "V6: intervention is sparse"

  -- confirm contraction
  assert plan.contraction_factor < 1.0
    emit "V6: gap contracts"

  report plan
}

-- V10: type safety (these must be rejected)
-- The following lines are commented because they MUST NOT compile:
-- let bad1 := 0.5#0         -- zero floor: ill-typed
-- close gap of r_raw        -- r_raw : Reading, not Cell
--   toward coherent using basis Glycolysis.enzymes
--   until coherent yield bad_plan
\end{lstlisting}

\subsection{Success and failure criteria}

\begin{criterion}[E5 success]
All ten validation steps pass on the five-node model. The type checker
rejects all three ill-typed programs. The gap-closure contraction
achieves $\mathrm{dist}(x_k, \mathcal{M}^*) < \flo$ within $k^* \leq 20$
iterations.
\end{criterion}

\begin{criterion}[E5 failure]
If any of steps V1--V9 fail on the five-node model, the CFC specification
has a bug that the constructive validation protocol was designed to catch.
If step V10 fails (the type checker accepts an ill-typed program), the
accountability type system is unsound.
\end{criterion}

% =====================================================================
\section{Implementation Plan}
\label{sec:implementation}
% =====================================================================

\subsection{Execution order}

The experiments are ordered by dependency and by diagnostic value:

\begin{enumerate}[leftmargin=1.4em,itemsep=3pt]
\item \textbf{E5} (CFC reference implementation) executes first, because
  every subsequent experiment runs on the CFC runtime and a broken
  runtime invalidates all downstream results. Estimated time: 1--2 weeks.
\item \textbf{E2} (spectral MHC binding) and \textbf{E3} (codon-usage
  overlap) execute in parallel, as both are FFT-and-dot-product
  computations on public sequence data with no shared dependencies.
  Estimated time: 2--3 days each.
\item \textbf{E1} (holonomy diagnosis) executes next, testing the core
  disease theory on real metabolic data. Estimated time: 1--2 weeks (data
  curation is the bottleneck).
\item \textbf{E4} (escape prediction) executes last, as it has the most
  complex data extraction (antibody CDR sequences from supplementary
  tables). Estimated time: 2--3 days.
\end{enumerate}

\subsection{Computational requirements}

No experiment requires GPU computation, distributed computing, or
proprietary software. Each runs on a single workstation with:
\begin{itemize}[leftmargin=1.4em,itemsep=1pt]
\item Python 3.10+ with NumPy, SciPy, scikit-learn.
\item The CFC reference interpreter (to be implemented as part of E5).
\item Internet access for database downloads (IEDB, Recon3D, NCBI,
  UniProt, IPD-IMGT/HLA).
\end{itemize}

Total estimated computation time across all five experiments:
$< 24$ CPU-hours.

\subsection{Reporting standard}

Each experiment will report:
\begin{enumerate}[leftmargin=1.4em,itemsep=2pt]
\item The raw data extraction script (reproducible download).
\item The CFC program (as given above, with any corrections from
  implementation).
\item The numerical results, keyed to the success/failure criteria.
\item Any corrections to the CFC syntax required by implementation
  (documenting the gap between specification and practice).
\item Any negative results, with the same prominence as positive ones.
\end{enumerate}

% =====================================================================
\section{Syntax Corrections Anticipated}
\label{sec:corrections}
% =====================================================================

The CFC programs above use constructs that extend the base grammar. We
anticipate three categories of correction during implementation:

\begin{enumerate}[leftmargin=1.4em,itemsep=3pt]
\item \textbf{Iteration.} The base CFC grammar has no \cf{foreach}
  construct. The \cf{scan} form of the \cf{spectrum} module subsumes
  iteration over database entries, but E1 and E5 require iteration over
  model parameters (rate constants, perturbation magnitudes). We
  anticipate adding:
  \begin{align*}
  \textit{foreachStmt} &\;::=\; \cf{foreach}\ \textit{ident}\
    \cf{in}\ \textit{expr}\ \cf{\{}\ \{\textit{stmt}\}\ \cf{\}}
  \end{align*}
  This is a sugar for sequential \cf{measure} calls and introduces no new
  semantic primitive.

\item \textbf{Collection.} The \cf{collect} clause of \cf{scan} returns a
  list. Operations on lists (mean, max, rank, sort, filter) are needed
  by every experiment. We anticipate a standard library of list
  operations, each typed as \cf{Reading -> Reading} (they transform
  readings into readings, preserving floor annotations).

\item \textbf{External data binding.} The \cf{import} statements above
  assume a binding between CFC identifiers and external databases (IEDB,
  Recon3D, GenBank). The CFC specification does not define this binding.
  We anticipate a \cf{data} module that maps CFC identifiers to file
  paths or URLs, with the contract that the imported data is read-only
  and carries no floor annotation (it is raw input, not a committed
  measurement).
\end{enumerate}

Each correction will be documented as a diff against the grammar of
\citep{sachikonye2026cfc}, \S9, and proved conservative by the argument
of \Cref{rem:conservative-circuit}: each new form binds an existing
semantic operation to a domain-specific type, advancing the
committed-measurement clock and preserving the accountability
invariants.

% =====================================================================
\section{Discussion}
\label{sec:discussion}
% =====================================================================

\subsection{What success would establish}

If all five experiments succeed by their stated criteria:

\begin{itemize}[leftmargin=1.4em,itemsep=2pt]
\item \textbf{E1} establishes that holonomy-based diagnosis works on real
  metabolic data, localises enzyme defects without being told which
  enzyme is deficient, and does so using a method (Wegscheider cycle
  consistency) that no existing metabolic analysis tool uses as a
  diagnostic classifier. This is a methods contribution to computational
  metabolomics.
\item \textbf{E2} establishes that a zero-training-data spectral method
  achieves competitive MHC binding prediction, with the distinctive
  property that its performance does not degrade on rare alleles. This is
  a methods contribution to computational immunology.
\item \textbf{E3} establishes a spectral signature of host tropism in
  protein sequences, supporting the housekeeping-mimicry and
  cross-kingdom non-overlap predictions. This is a novel finding in
  computational virology.
\item \textbf{E4} establishes that the spectral escape landscape captures
  real immunological selection pressure, ranking known escape mutations
  in the top decile of candidates. This has direct implications for
  vaccine design and pandemic preparedness.
\item \textbf{E5} establishes that CFC is implementable as specified and
  that its accountability guarantees (no zero-residue value, no blind
  therapy) hold on a concrete model. This is a software contribution.
\end{itemize}

\subsection{What failure would establish}

Each failure criterion is stated so that a negative result is informative
rather than merely disappointing:

\begin{itemize}[leftmargin=1.4em,itemsep=2pt]
\item \textbf{E1 failure} distinguishes data-quality limitations (healthy
  circuit fails KCL) from theory limitations (healthy circuit passes but
  disease circuit shows no holonomy defect).
\item \textbf{E2 failure} at AUC $< 0.70$ would show that Fourier
  magnitudes without phase are insufficient for binding prediction,
  motivating complex-valued embeddings.
\item \textbf{E3 failure} would show the spectral embedding does not
  capture host-specificity, motivating codon-level rather than
  amino-acid-level analysis.
\item \textbf{E4 failure} would show the spectral escape landscape does
  not capture real selection, motivating structural rather than
  sequence-based escape prediction.
\item \textbf{E5 failure} would identify a bug in the CFC specification.
\end{itemize}

\subsection{What is not tested}

These experiments do not test:
\begin{itemize}[leftmargin=1.4em,itemsep=2pt]
\item The template-matching epidemiology of \citep{sachikonye2026template}
  (requires population-level data not currently in scope).
\item The attention-allocated retrieval of
  \citep{sachikonye2026attention} (requires live federated endpoints).
\item The four-column alignment of \citep{sachikonye2026fourcolumn}
  (requires a canonical response construction, which that paper
  identifies as an open problem).
\item Any wet-lab experiment.
\end{itemize}

% =====================================================================
\section{Conclusion}
\label{sec:conclusion}
% =====================================================================

We have designed five experiments that test the principal claims of the
circuit-based disease framework against public data, using the CFC
domain-specific language as the implementation vehicle. Three syntax
extensions to CFC are specified: a circuit module, a spectrum module, and
a network module, each conservative in the sense that it binds existing
CFC semantic operations to domain-specific types without introducing new
reduction rules.

Each experiment has a quantitative success criterion and a quantitative
failure criterion stated in advance, so that the outcome---positive or
negative---is informative. No experiment requires wet-lab work,
proprietary data, or computational resources beyond a single workstation.
The total estimated computation time is under 24 CPU-hours.

The experiments are executable as written. What remains is to execute
them.

\bibliographystyle{plainnat}
\bibliography{references}

\end{document}